\chapter{Introduction}

With technical progress in the field of robotics and related areas, autonomous robots are now able to execute many tasks in factories, even without modifications to the environment like painted lines. However, robots struggle to conquer everyday human environments. Domestic environments might seem very chaotic at the first glance, with a wide variety of different types of environments and huge differences in appearance. However, those environments have an underlying structure, which is used by humans all the time to carry out tasks efficiently. Indoor environments are usually segmented into rooms, which are separated by walls and only connected by doors. A room consists of one or a small number of areas, where each area has specific purposes, e.g. a kitchen is used for baking and cooking. For each purpose specific objects are used, e.g. a pan for cooking, which can be assumed to be somewhere around. 

The goal of this thesis is to equip a robot with such commonsense knowledge about indoor environments to execute a task efficiently. The investigated task is to search an object in an unknown open indoor environment. This is a relative simple task, but can benefit a lot from commonsense knowledge about indoor environments. It is also an essential skill for a robot working in domestic environments and therefore research on this topic will be useful in future robotic systems. During the execution of this task the robot has to find the best next view pose. In this thesis the best next view pose is the pose that minimizes the expected search time. The robot has to decide which areas are promising to contain the searched object. It also has to find a balance between exploration and the search of already explored areas because the environment is unknown at the start of a search. Here the commonsense knowledge is important to make better estimations about not or only vaguely seen areas. 

The system presented in this work uses a hierarchical approach. The environment is seen as a topology of rooms at a higher abstraction level and the details are handled for every room separately. Within a room convolutional neural networks (CNNs) for object detection and place classification are used to create semantic maps containing information about the positions where objects were seen and the room types of areas in the room. Using the results of the object detection, the place classification and the connection between the type of an area and the objects usually located in that area, object probabilities are estimated, even for not or only vaguely seen areas. To achieve this a new mapping system was developed, capable to create hybrid maps, which combine the room topology and semantic maps of the rooms. The developed mapping system is also able to cope with the uncertain results of the sensors and the CNNs. The hybrid map is then used by the developed hierarchical search planner, which first selects the next room to search and then plans the search within that room, to execute the search task efficiently. 

Similar systems were already developed, e.g. by Aydemir et al. \cite{main}. However, the generality of the used CNNs are a step towards the flexibility, which will be necessary in future robots used for example in households. The proposed system was tested in common indoor environments, where the behavior was investigated and a comparison with a system using no commonsense knowledge was done.